% Verbatim copy of /Users/alingalatan/workspaces/grothe/handoff/research-degree5/external-2026-11-15/projection-attack-package/proofs.tex (file dated 2026-11-15). Only this line was added.
% Intended filepath: grothendieck_projection_attack/proofs.tex
\documentclass[11pt]{article}
\usepackage[a4paper,margin=27mm]{geometry}
\usepackage{amsmath,amssymb,amsthm,mathtools,hyperref}
\newtheorem{theorem}{Theorem}
\newtheorem{proposition}[theorem]{Proposition}
\newtheorem{lemma}[theorem]{Lemma}
\newcommand{\E}{\mathbb E}
\newcommand{\C}{\mathbb C}
\newcommand{\R}{\mathbb R}
\newcommand{\ph}{\operatorname{ph}}
\newcommand{\Rea}{\operatorname{Re}}
\newcommand{\HH}{\mathcal H}
\title{Retaining the first-chaos penalty\\\large Rigorous partial progress toward a lower bound above 1.402}
\author{Research calculations}
\date{14 November 2026}
\begin{document}
\maketitle
\begin{abstract}
This note does not establish a new numerical bound for the complex
Grothendieck constant. It proves an exact support-function reduction that
retains the first-chaos penalty, a global improvement of the remote row,
and the resulting local two-sector cost matrix. It also supplies a certified
counterexample to an auxiliary weighted-contraction positivity claim.
The missing step is uniform finite-amplitude control at general remote bases.
\end{abstract}

\section{Normalization and the retained penalty}
All Gaussian coordinates are standard complex Gaussians, so
$\E|Z|^2=1$. Inner products are $\langle f,g\rangle=\E\bar f g$.
Let $P$ and $Q$ be orthogonal projections onto $\HH_{1,0}$ and
$\HH_{2,1}$, respectively. Put
\[
 \epsilon=\frac{1891}{2500},\quad p=\epsilon^{-1},\quad
 \ell=\frac{1057}{1250},\quad b=\frac{179}{1250},\quad
 a=\ell-2b=\frac{699}{1250},\quad L=Z(a+b|Z|^2).
\]
The analytic statements involving $L$ below also hold for any $a,b>0$.
Define
\[
 \mathcal F_L(W)=\E(|L+W|+|L-W|),\qquad
 \Lambda(G)=\mathcal F_L(G)-p\|G\|_2^2,
\]
and the smaller functional
\begin{equation}\label{eq:retained}
 \widehat\Lambda(G)=
 \min_{u\in\HH_{1,0}}
 \{\|u\|_2^2+\mathcal F_L(G-u)\}-p\|G\|_2^2.
\end{equation}
The penalty is on the original cubic $G$, not on $G-u$.
In a finite Gaussian dimension the minimum exists and is unique: its objective
is continuous, coercive, and strongly convex in the finite-dimensional variable $u$.
In particular $\widehat\Lambda(G)\le\Lambda(G)$.

\begin{theorem}[Exact penalized support duality]\label{thm:duality}
For fixed $L,G\in L^2$,
\begin{align*}
 &\sup_{|f|,|g|\le1}\left\{
 \Rea\E[\bar f(L+G)+\bar g(L-G)]
 -\left\|P\frac{f-g}{2}\right\|_2^2\right\}\\
 &\hspace{25mm}=\min_{u\in\HH_{1,0}}
 \{\|u\|_2^2+\mathcal F_L(G-u)\}.
\end{align*}
Consequently, for unimodular responses $f,g$, put $h=(f+g)/2$,
$k=(f-g)/2$. Rotate coordinates so $Ph=tZ$, $t\ge0$, and put
$G=\epsilon Qk$. Then
\begin{equation}\label{eq:row}
 \Rea\langle f,(P-\epsilon Q)g\rangle
 \le t^2-2\ell t+\frac{2b^2}{\epsilon}+\widehat\Lambda(G).
\end{equation}
\end{theorem}
\begin{proof}
Use the identity
\[
 -\|Pk\|_2^2
 =\min_{u\in\HH_{1,0}}\{\|u\|_2^2-2\Rea\langle u,k\rangle\}.
\]
The sets of disk-valued responses are weakly compact convex subsets of $L^2$.
The displayed expression is affine and weakly continuous in $(f,g)$ for fixed
$u$, and convex continuous in $u$. Convex minimax therefore exchanges the
supremum and infimum. The pointwise support of the unit disk is the modulus,
which gives the claimed formula. For the upper-bound direction alone, no
minimax theorem is needed: complete the square and take the support bound
for each fixed $u$, then infimize.

For the row inequality, orthogonality gives
\[
 \Rea\langle f,(P-\epsilon Q)g\rangle
 =t^2-\|Pk\|_2^2-\epsilon\|Qh\|_2^2+\epsilon\|Qk\|_2^2.
\]
The Wick polynomial $Z(|Z|^2-2)$ has squared norm $2$. Completing the
$Qh$ square and using $G=\epsilon Qk$ bounds the right side by
\[
 t^2-2\ell t+2b^2/\epsilon
 +2\Rea\langle L,h\rangle+2\Rea\langle G,k\rangle
 -\|Pk\|_2^2-p\|G\|_2^2.
\]
Apply the penalized support formula.
\end{proof}

This is a relaxation of the actual-response problem: the support supremum
does not retain $G=\epsilon Qk$ or $Ph=tZ$ as constraints. It nevertheless
retains more information than the original $\Lambda$ row. A uniform cap on
$\widehat\Lambda$ would suffice in the same row envelope. Nothing here supplies
that cap.

\section{A global finite-amplitude improvement for remote cubics}
Let $q$ depend only on Gaussian coordinates independent of $Z$. Set
\[
 M(w)=\E_Z|L+w|,\qquad m(w)=\E_Z\ph(L+w),\qquad
 j_0=\frac12\E\frac1{|L|}.
\]
Here $\ph(0)$ may be set to zero; the exceptional sets do not affect the integrals.
Let $P_{\rm rem}$ denote the projection onto remote complex linear functions.

\begin{lemma}[Gaussian-mixture representation and smoothness]\label{lem:smooth}
The law of $L$ is a mixture of centered proper complex Gaussian laws.
The gradient $m=\nabla M$, interpreted on $\R^2$, is globally $j_0$-Lipschitz.
\end{lemma}
\begin{proof}
Write $s=|L|^2=t(a+bt)^2$ and let $t=t(s)$ be the increasing inverse.
The planar density at squared radius $s$ is
\[
 \rho(s)=\frac{e^{-t(s)}}{\pi(a+bt(s))(a+3bt(s))}.
\]
Put $v(t)=((a+bt)(a+3bt))^{-1}$, so $t'(s)=v(t(s))$.
Both $v$ and $H(t)=e^{-t}v(t)/\pi$ are completely monotone.
If $H$ is completely monotone, so is $-vH'$, by closure under products.
Iteration of this observation shows
\[
 (-1)^n\rho^{(n)}(s)=(-v\partial_t)^nH(t(s))\ge0.
\]
The Laplace representation of a completely monotone function gives
$\rho(s)=\int e^{-\tau s}\nu(d\tau)$ with $\nu\ge0$.
Normalization gives $\int(\pi/\tau)\nu(d\tau)=1$ and no atom at $\tau=0$.
Thus this is a mixture of the planar densities
$(\tau/\pi)e^{-\tau|w|^2}$.

For a proper Gaussian $N$ of variance $v$, write
$\E\ph(N+w)=w j_v(|w|^2)$. The Laplace formula for the square root yields
\[
 j_v(s)=\frac1{\sqrt\pi}\int_0^\infty
 \frac{t^{-1/2}}{(1+vt)^2}
 \exp\left(-\frac{st}{1+vt}\right)dt,
 \qquad j_v(0)=\frac{\sqrt\pi}{2\sqrt v}.
\]
The two real Jacobian eigenvalues are $j_v(s)$ and
$j_v(s)+2s j_v'(s)$. Both have absolute value at most $j_v(0)$, since
$e^{-x}\le1$ and $|(1-2x)e^{-x}|\le1$ for $x\ge0$.
Integrating over the Gaussian mixture gives the claimed Lipschitz constant
$\E[\sqrt\pi/(2\sqrt v)]=\frac12\E|L|^{-1}=j_0$.
\end{proof}

\begin{theorem}[Global retained-projection improvement]\label{thm:remote}
For every remote $q\in L^2$,
\begin{equation}\label{eq:globalgain}
 \boxed{\quad
 \widehat\Lambda(q)\le\Lambda(q)
 -\frac{\|P_{\rm rem}m(q)\|_2^2}{1+j_0}.
 \quad}
\end{equation}
In particular the inequality holds at every amplitude for every remote cubic,
in every finite Gaussian dimension.
\end{theorem}
\begin{proof}
Circular symmetry gives $\mathcal F_L(q)=2\E M(q)$.
Lemma~\ref{lem:smooth} gives, for every remote linear $u$,
\[
 \mathcal F_L(q-u)\le\mathcal F_L(q)
 -2\Rea\langle u,m(q)\rangle+j_0\|u\|_2^2.
\]
Use this in \eqref{eq:retained} and choose
$u=P_{\rm rem}m(q)/(1+j_0)$. No cubic moment estimate or independence between
$q$ and the remote linear functions is used.
\end{proof}

The correction is nonlinear in $q$. It does not vanish merely because the
leading contraction $P_{\rm rem}(q|q|^2)$ vanishes. Conversely, this theorem
alone does not handle the coupled $A$ and $B$ sectors.

\section{The exact corrected local A/B cost matrix}
Let $r_A,r_B$ be remote complex linear functions, with arbitrary mutual
covariance, and write $T=|Z|^2$ and
\[
 H=(T-1)r_A+\frac{Z^2\overline{r_B}}{\sqrt2}.
\]
Define three scalar integrals
\[
 I_k=\int_0^\infty\frac{t^{k-1/2}e^{-t}}{a+bt}\,dt,
 \qquad k=0,1,2,
\]
and
\begin{align*}
 j_0&=I_0/2,&g_0&=(I_1-I_0)/2,&d_0&=-I_1/(2\sqrt2),\\
 k_A&=p-(I_2-2I_1+I_0)/2,&
 k_B&=p-I_2/4,&m_0&=(I_2-I_1)/(2\sqrt2).
\end{align*}

\begin{theorem}[Local Schur-complement gain]\label{thm:matrix}
As the real number $s\to0$,
\[
 \widehat\Lambda(sH)=2\E|L|-s^2\mathcal Q(r_A,r_B)+o(s^2),
\]
where
\begin{align*}
 \mathcal Q(r_A,r_B)
 &=k_A\|r_A\|_2^2+k_B\|r_B\|_2^2
   +2m_0\Rea\langle r_A,r_B\rangle\\
 &\quad+\frac{\|g_0r_A+d_0r_B\|_2^2}{1+j_0}.
\end{align*}
At the peak row, its coefficient matrix is enclosed by the independently
computed Arb values
\[
 \begin{pmatrix}
 0.60683885483874355964\ldots&0.26291259657063581075\ldots\\
 0.26291259657063581075\ldots&1.01343807031137005476\ldots
 \end{pmatrix},
\]
and has smallest eigenvalue greater than $0.477$.
\end{theorem}
\begin{proof}
For $x\ne0$, the quadratic term of $|x+y|$ is
\[
 \frac1{2|x|}\left(|y|^2-
       \frac{\Rea(\bar x y)^2}{|x|^2}\right).
\]
Apply this to $L\pm sH$ and average over the phase of $Z$ and the proper
Gaussian remote linear forms. The coefficient in $\Lambda(sH)$ is
\[
 -k_A\|r_A\|_2^2-k_B\|r_B\|_2^2
 -2m_0\Rea\langle r_A,r_B\rangle.
\]
For completeness, the coefficient in $\mathcal F_L(sH-su)$, where $u$ is
remote linear, is
\begin{align*}
 &\frac{I_2-2I_1+I_0}{2}\|r_A\|_2^2
 +\frac{I_2}{4}\|r_B\|_2^2
 -\frac{I_2-I_1}{\sqrt2}\Rea\langle r_A,r_B\rangle\\
 &\quad+j_0\|u\|_2^2
 -2\Rea\langle g_0r_A+d_0r_B,u\rangle.
\end{align*}
Adding $\|su\|_2^2$ and minimizing gives
$u=(g_0r_A+d_0r_B)/(1+j_0)$ and subtracts the squared-norm term in the theorem.
The visible linear direction is uncoupled at this order, and its minimizing
coefficient is zero.

These Taylor expansions are legitimate despite $L(0)=0$: their remainders
are dominated by a constant times the squared perturbation divided by $|L|$.
This is integrable for the finitely many linear and cubic directions involved,
because $a>0$ and the radial singularity is only $t^{-1/2}$.
Strong convexity of the correction-variable objective implies that its
minimizer is $O(s)$, so its minimized quadratic expansion is the Schur complement.

Finally,
\[
 I_0=\frac{\pi}{\sqrt{ab}}e^{a/b}\operatorname{erfc}\sqrt{a/b},\quad
 I_1=\frac{\sqrt\pi-aI_0}{b},\quad
 I_2=\frac{\sqrt\pi/2-aI_1}{b}.
\]
The script evaluates these expressions using exact rational inputs and
160-bit Arb ball arithmetic and checks the stated rational inequalities.
\end{proof}

\textbf{Scope.} This is a local expansion at $G=0$. It is not a finite-amplitude
bound and is not a Hessian bound at an arbitrary remote cubic $q$.
Before eliminating the projection variable the local A coefficient is
$0.320054503608992\ldots$; afterward it is $0.606838854838743\ldots$.
The latter must not be inserted as a global cost constant in a lower-bound certificate.

\section{A necessary maximizer condition and a failed positivity route}
There is a useful constraint on a maximizer that is stable under adding
fresh Gaussian coordinates. If $Y$ is a fresh standard complex Gaussian and
$P_2\in\HH_{1,1}$ uses only the old coordinates, then $YP_2\in\HH_{2,1}$.
At a local maximum $G$, assuming the indicated second derivative is integrable,
\begin{equation}\label{eq:stability}
 \E\left[\frac{|P_2|^2}{4}
 \left(\frac1{|L+G|}+\frac1{|L-G|}\right)\right]
 \le p\E|P_2|^2.
\end{equation}
Indeed, conditionally on the old coordinates $YP_2$ is proper Gaussian, so
the second derivative of $\Lambda(G+sYP_2)$ is twice the left side minus
$2p\E|P_2|^2$. The same argument applies to $\bar YP_2$ with
$P_2\in\HH_{2,0}$. This does not assert that a maximizer exists in the
dimension-uniform problem, nor that these necessary conditions suffice.

One tempting way to use such stability conditions is to assert positivity
of weighted cubic contractions. The following explicit example rules out
one such assertion.

\begin{proposition}[A weighted contraction can be negative]\label{prop:counter}
Let $X,Y$ be independent standard complex Gaussians and
\[
 q=\frac{X\,[5(|X|^2-2)+6i(|Y|^2-1)]}{\sqrt{86}}.
\]
Then $q\in\HH_{2,1}$, $\E|q|^2=1$, and, with Wirtinger derivatives,
\[
 \boxed{\quad
 \Rea\E\left[e^{-|q|^2}\bar q^{\,2}
 \sum_{j=X,Y}(\partial_jq)(\partial_{\bar j}q)\right]<-\frac1{50}.
 \quad}
\]
\end{proposition}
\begin{proof}
Put $R=|X|^2$, $S=|Y|^2$, $V=S-1$, $c=25/86$, $d=36/86$.
The variance is $2c+d=1$. Direct differentiation gives
$|q|^2=R[c(R-2)^2+dV^2]$ and
\begin{align*}
 \Rea\left[\bar q^2\sum_jq_jq_{\bar j}\right]
 =R^2\{&c(R-2)^2[c(2R-2)-d]-cd(R-2)^2V\\
       &+(d^2-2cd)V^2+d^2V^3\}.
\end{align*}
For $v>0$ define $M_k(v)=\int_0^\infty(s-1)^k e^{-s-v(s-1)^2}\,ds$.
Completing a square and then integrating by parts gives
\begin{align*}
 M_0(v)&=\frac{\sqrt\pi}{2\sqrt v}
 e^{-1+1/(4v)}\operatorname{erfc}\frac{1-2v}{2\sqrt v},\\
 M_1&=\frac{e^{-v}-M_0}{2v},\quad
 M_2=\frac{M_0-M_1-e^{-v}}{2v},\quad
 M_3=\frac{2M_1-M_2+e^{-v}}{2v}.
\end{align*}
Thus the claimed expectation is the integral on $r\ge0$ of
\begin{align*}
 f(r)=r^2e^{-r-cr(r-2)^2}\{&c(r-2)^2[c(2r-2)-d]M_0(dr)\\
 &-cd(r-2)^2M_1(dr)+(d^2-2cd)M_2(dr)+d^2M_3(dr)\}.
\end{align*}
The certificate uses three intervals. On $[0,1/8]$,
$|M_k|\le\E(S+1)^k$ gives the bounds $1,2,5,16$ for $k=0,1,2,3$.
Consequently $|f(r)|\le6r^2$: the rational coefficient bound checked by the
script is
\[
 4c(2c+d)+8cd+5|d^2-2cd|+16d^2<6.
\]
This interval contributes at most $1/256$.
For $r\ge8$, the same moment bounds give
$|f(r)|\le10r^5e^{-r-288c}$, whose integral is
\[
 10e^{-288c-8}\,120\sum_{j=0}^5\frac{8^j}{j!}<10^{-30}.
\]
On $[1/8,8]$, the script sums interval enclosures of the whole integrand
over dyadic cells of width $2^{-12}$. It adds both preceding upper bounds
and verifies that the resulting upper endpoint is below $-1/50$.
All inputs and cell endpoints are exact rationals. This is an inclusion
calculation, not a floating-point quadrature error estimate.
\end{proof}

This proposition is \emph{not} a counterexample to Conjecture K. It prevents
using a pointwise-in-scale positivity assertion for this contraction in a proof.

\section{What remains before a new endpoint}
To prove $K_G^{\C}>1.402$ using this approach, the completed row envelope must
have a certified cap below
\[
 \frac{500}{701}=0.71326676176890156918\ldots.
\]
The analogous threshold for $1.4022$ is $5000/7011$.
The retained-penalty duality is exact; the remote improvement is global;
the two-sector matrix is only local. A proof still needs a uniform estimate
for \eqref{eq:retained} with general coupled remote and visible sectors, and
then the row-uniform numerical certificate. The remote result and local
matrix do not supply those missing estimates. No new upper bound is claimed.

\paragraph{Verification record.}
Run \texttt{python verify.py} with \texttt{python-flint} installed.
The included \texttt{results.txt} is the output of the executed verifier.
The analytic arguments in this note are written proofs; they have not been
formalized in Lean. No changes were made to the remote repository.
\end{document}
